\section{Theoretical Results}

In this section, we answer the questions posed in the previous section. We first derive the optimal detector for the single-step decision problem and then extend this analysis to the sequential setting to characterize the fundamental limit of sample efficiency. Let $n$ denote the cardinality of the sample space $\cV$ and $\mathbf{e}_v$ denote the one-hot vector (i.e., Dirac measure) that assigns mass $1$ to $v \in \cV$.

\subsection{Optimal log-growth rate}

We begin by solving the robust log-optimality problem defined in \Cref{prob:robust-log-optimal}. The following theorem provides a closed-form solution for both the optimal e-value and the worst-case-optimal generator behavior.

\begin{theorem}[Log-growth rate]\label{thm:log-growth}
The optimal e-value that solves the objective in Eq.~\eqref{eq:log-growth} is given by:
\begin{align}\label{eq:optimal-e}
e^*(v,s)
=
\begin{cases}
\frac{1-\delta/2}{p_0(s)}, & s=v,\\
\frac{\delta}{2(n-1)p_0(s)}, & s\neq v,
\end{cases}
\end{align}
where $n = |\cV|$. Furthermore, for any target distribution $q \in \cQ(p_0,\delta)$ decomposed as $q = \sum_{i=1}^k \lambda_i q_i$ where $q_i = p_0 + \frac{\delta}{2}\cdot (\mathbf{e}_{v_i} - \mathbf{e}_{s_i}), v_i \neq s_i \in \cV$ are the extreme points of $\cQ(p_0,\delta)$, 
the optimizer of the inner maximization problem (the generator's optimal coupling) is:
\begin{align}\label{eq:coupling}
    w^* = \sum_{i=1}^k \lambda_i \cdot \left(\frac{\delta}{2} \cdot (\mathbf{e}_{v_i}\mathbf{e}_{s_i}^\top - \mathbf{e}_{s_i}\mathbf{e}_{s_i}^\top) + \diag(p_0)\right).
\end{align}
The optimal robust log-growth rate is equal to:
\begin{align}\label{eq:optimal-log-growth}
    J^* = h + \left(1-\frac{\delta}{2}\right)\log\left(1-\frac{\delta}{2}\right)
 + \frac{\delta}{2}\log\left(\frac{\delta}{2(n-1)}\right),
\end{align}
where the term $h = H(p_0)$ is the Shannon entropy of the anchor distribution $p_0$.
\end{theorem}

The result in \Cref{eq:optimal-log-growth} fundamentally express the optimal growth rate of e-values in Anchored E-Watermarking. The first term is the entropy of the anchor distribution, which mathematically formalizes the intuition that watermarking is inherently easier for distributions with high entropy. 
The second term can be written as $-H(\nu_\delta)$, the negative entropy of the categorical distribution $\nu_\delta = (1-\delta/2, \delta/(2(n-1)),\dots,\delta/(2(n-1))) \in \Delta([n])$, which is a decreasing function of $\delta$. This explicitly characterizes the power-robustness tradeoff determined by the tolerance parameter $\delta$. 
Written as $H(p_0) - H(\nu_\delta)$, the optimal log-growth rate is the achieved information rate of an Additive Noise Channel $Y = X \oplus Z$ where $X \sim p_0, Z \sim \nu_\delta$. 

\begin{remark}[Relationship to SEAL~\citep{huang2025watermarking}]
SEAL proposes to use a smaller language model as $p_0$ and speculative decoding as the generator.
Note that the optimizer of the inner problem given in Eq.~\eqref{eq:coupling} is exactly the maximal coupling given by speculative decoding. Therefore, \Cref{thm:log-growth} confirms that the choice of generator in \citet{huang2025watermarking} is optimal. However, the optimal detector given by Eq.~\eqref{eq:optimal-e} is different from the detection rule in SEAL: this gap is corroborated by experiments in Section~\ref{sec:exp-real}.
\end{remark}

\subsection{Sample efficiency}

Having characterized the one-step optimal growth rate, we now analyze the long-term performance of the watermark in a sequential setting. The following theorem connects the log-growth rate $J^*$ to the sample complexity required to reject the null hypothesis against an adaptive adversary.

\begin{theorem}[Expected stopping time]
\label{thm:stopping-time}
Let $\mu(\cA,\cG)$ be the stochastic process defined in \Cref{prob:stopping-time} and $\tau_\alpha(e)$ be the stopping time under e-value $e$. Let $J^*$ be the optimal rate defined in Eq.~\eqref{eq:optimal-log-growth}. We have the following bounds on sample efficiency:
\begin{itemize}
    \item \textbf{Lower bound (converse):} For any valid e-value $e \in \cE$:
\begin{align*}
    \sup_{\cA} \inf_{\cG} \liminf_{\alpha \to 0} \frac{\mathbb{E}_{\mu(\cA,\cG)}[\tau_\alpha(e)]}{\log(1/\alpha)} \geq \frac{1}{J^*}.
\end{align*}
\item \textbf{Upper bound (achievability):} For the optimal e-value $e^*$ defined in \Cref{thm:log-growth}:
\begin{align*}
    \sup_{\cA} \inf_{\cG} \liminf_{\alpha \to 0} \frac{\mathbb{E}_{\mu(\cA,\cG)}[\tau_\alpha(e^*)]}{\log(1/\alpha)} = \frac{1}{J^*}.
\end{align*}
\end{itemize}
\end{theorem}


\Cref{thm:stopping-time} establishes $1/J^*$ as the fundamental information-theoretic limit of sample complexity for Anchored E-Watermarking. It demonstrates that the e-value $e^*$ derived from the one-step greedy optimization is not only locally optimal but also globally optimal for sequential testing. Crucially, this optimality holds even against an adaptive adversary that can vary the target distribution at every step, provided the distribution remains within the $\delta$-neighborhood of the anchor. 

\begin{remark}
[Relationship with the rates in \citet{huang2023towards}]
\citet{huang2023towards} shows that the minimum number of samples required to watermark with Type I error $\alpha$ scales as $\log(1/\alpha)\cdot\frac{\log(1/h)}{h}$, where $h$ is the average entropy per token. While our rate $\frac{\log(1/\alpha)}{h}$ has the same scaling in terms of $\alpha$, its dependence on the entropy $h$ is improved. This is because \citet{huang2023towards} study an asymptotic regime where $h \to 0$, while we fix an anchor distribution with lower bounded entropy arising from the condition $\inf_{v \in \cV}p(v) > \delta$.
\end{remark}